\chapter{Algorithmic Market Making \& Adverse Selection}
\label{ch:algorithmic_market_making}

If AMPS relied entirely on generative Large Language Models (LLMs) to provide resting liquidity, the simulated market would instantly collapse. The inherent latency of LLM inference -- even when heavily quantized -- is fundamentally incompatible with the continuous, microsecond-level quoting obligations required to maintain a functional spread. Therefore, the simulation architecture mandates the deployment of deterministic \textbf{Algorithmic Market Makers (AMMs)}. This chapter explores the mathematical logic driving these AMMs and how their defensive behaviors act as the primary catalyst for flash crashes.

\section{The Function of the Market Maker}

A market maker is an algorithmic agent designed to provide continuous liquidity to the market. It does this by simultaneously quoting both a Bid and an Ask price. The market maker does not hold a directional opinion on the true value of the asset; instead, its goal is to earn the \textit{Spread} (the difference between its Ask and its Bid) over thousands of round-trip trades.

If the market maker buys at $\$99.99$ and immediately sells at $\$100.01$, it captures $\$0.02$ in risk-free profit. However, this theoretical profit is constantly threatened by two primary risks: \textbf{Inventory Risk} and \textbf{Adverse Selection}.

\section{Inventory Risk and Mean Reversion}

An AMM must maintain a balanced inventory. If the market experiences a temporary influx of retail sellers, the AMM will continuously buy the asset, accumulating a massive long position. If the price permanently drops, the AMM suffers catastrophic capital losses.

To mitigate this, AMMs implement \textit{Inventory Control}. As their inventory of an asset grows linearly, they skew their quotes exponentially to attract the opposite order flow. 
For example, if an AMM is dangerously long (holding too much of the asset), it will:
\begin{itemize}
    \item Lower its Bid price (to discourage further buying).
    \item Lower its Ask price (to aggressively attract sellers and offload its inventory).
\end{itemize}

While inventory control is critical for normal market equilibrium, it breaks down entirely during a macroeconomic panic, leading us to the concept of Adverse Selection.

\section{Adverse Selection and Toxic Flow}

Adverse Selection occurs when a market maker trades against an informed counterparty. In the context of AMPS, when the Macro God Agent broadcasts a severe negative shock, the ``Whale'' agents receive this information instantly. They generate \textbf{Toxic Flow} -- aggressive market sell orders driven by deterministic knowledge that the asset is overvalued \cite{glosten1985bid}.

\subsection{The Glosten-Milgrom Formulation}

The classic Glosten-Milgrom model \cite{glosten1985bid} provides the mathematical framework for understanding how AMMs react to toxic flow. Let $V \in \{V_{low}, V_{high}\}$ be the true fundamental value of the asset, which is known to informed traders but unknown to the AMM. The AMM only observes the order flow.

When the AMM observes an overwhelming ratio of sell orders to buy orders (an extreme Order Flow Imbalance, or OFI), its internal probability estimation that $V = V_{low}$ approaches $1.0$.

To protect its capital, the AMM mathematically \textit{must} widen its spread $S$. The formula for the required spread to break even against informed traders is directly proportional to the probability of trading with an informed trader ($\mu$):

\begin{equation}
    S \propto \frac{\mu}{1 - \mu} \cdot \sigma_V
\end{equation}

where $\sigma_V$ represents the volatility of the asset's true value. As $\mu \to 1$ (representing a pure panic where all incoming flow is toxic), the required spread $S \to \infty$. 

\begin{figure}[htbp]
\centering
\begin{tikzpicture}[
    scale=0.9, transform shape,
    box/.style={draw, rectangle, minimum width=3cm, minimum height=0.8cm, align=center},
    arr/.style={->, thick, >=stealth}
]
    \node[box, fill=red!10] (toxic) {Toxic Order Flow\\(Whale Panic Selling)};
    \node[box, fill=orange!10] (ofi) [below=0.8cm of toxic] {Severe Order Flow\\Imbalance (OFI) Detected};
    \node[box, fill=yellow!10] (bayes) [below=0.8cm of ofi] {AMM Bayesian Update\\($P(V_{low}) \to 1.0$)};
    \node[box, fill=blue!10] (withdraw) [below=0.8cm of bayes] {Exponential Spread\\Widening / Liquidity Void};
    
    \draw[arr] (toxic) -- (ofi);
    \draw[arr] (ofi) -- (bayes);
    \draw[arr] (bayes) -- (withdraw);
\end{tikzpicture}
\caption{The deterministic logic pipeline of an Algorithmic Market Maker under toxic flow conditions.}
\label{fig:amm_logic}
\end{figure}

\section{Limit Up -- Limit Down (LULD) Circuit Breakers}

Because algorithms will perfectly execute Equation 3.1, a severe panic will cause the AMMs to instantly cancel all resting bids, creating a \textit{Liquidity Void}. If a generative retail agent submits a market sell order into a liquidity void, the order will execute at a price of $\$0.00$, destroying the market simulation.

To prevent this, traditional finance (TradFi) exchanges implement regulatory \textbf{Circuit Breakers}, such as the Limit Up -- Limit Down (LULD) mechanism \cite{sec2012luld}. LULD explicitly pauses trading for 5 minutes if the price of an asset moves beyond a specified percentage band (e.g., 5\% or 10\%) within a rolling 5-minute window. 

The inclusion of LULD dynamics is what strictly separates AMPS from unregulated cryptocurrency simulations, forcing the generative agents to comprehend regulatory halts and the physics of unexecutable panic.
